\section{Research Project Experience}

\setlength{\arrayrulewidth}{.01pt}
\arrayrulecolor[HTML]{C1C1C1}
\setlength{\tabcolsep}{8pt}

\noindent
\begin{tabular}{r|L{14cm}}

    & \textbf{Selected NASA DEVELOP Projects} {\hfill| \small Advisor: Jonathan \textsc{Horton}, Ph.D.}\\ 
    \textsc{2023} 
    & \footnotesize{\emph{Pacific Northwest Health and Air Quality}}\\
    & \footnotesize{\textsc{\underline{Role:} Project Manager and Supervisor}} \\
    & \footnotesize{P}
    \\&\\
    
    \textsc{2023} 
    & \footnotesize{\emph{Northeast Alaska Climate}}\\
    & \footnotesize{\textsc{\underline{Role:} Project Manager and Supervisor}} \\
    & \footnotesize{P}
    \\&\\
    
    \textsc{2022} 
    & \footnotesize{\emph{Gatlinburg and Beatty Wildfires}}\\
    & \footnotesize{\textsc{\underline{Role:} Project Manager and Supervisor}} \\
    & \footnotesize{P}
    \\&\\
    
    \textsc{2022} 
    & \footnotesize{\emph{Vermont and New Hampshire Ecological Forecasting}: Monitoring Trends in Tree Defoliation Due to \textit{Lymantria dispar} Outbreaks to Predict Future Hardwood Tree Mortality and Health Impacts}\\
    & \footnotesize{\textsc{\underline{Role:} Research Associate}} \\
    & \footnotesize{\textsc{\underline{Co-Investigators:} Morgan Dean, Jacob Orser, Seamore Zhu}}\\
    & \footnotesize{Modeled tree defoliation due to the invasive month \textit{Lymantria dispar} using Landsat 7/8, Sentinel-2, SMAP, and Terra MODIS satellite data to support forest management efforts by the Forest Ecosystem Monitoring Cooperative, the Vermont Agency of Agriculture, Food and Markets, the University of New Hampshire Cooperative Extension, and the New Hampshire Division of Forests and Lands, Forest Health Program
    }
    \\&\\
    
    \textsc{2021} 
    & \footnotesize{\emph{Carolina Coastal Plain Ecological Forecasting}: Utilizing NASA Earth Observations to Map Suitable Venus Flytrap Habitat to Inform Conservation Efforts, Seed Banking, and Reintroduction in the Carolina Coastal Plain and Sandhills Region}\\
    & \footnotesize{\textsc{\underline{Role:} Research Associate}} \\
    & \footnotesize{\textsc{\underline{Co-Investigators:} Jayne Lampley, Monika Rock, Ashna Siddhi, Seamore Zhu}}\\
    & \footnotesize{Modeled present-day and future suitable Venus flytrap habitat using all high-confidence population observations recorded since 1980, 23 environmental predictor variables, and land-use change variables to inform seed banking and reintroduction}
    \\&\\
    
    & \textbf{Selected Projects at the University of North Carolina Asheville} \\ 
    \textsc{2019} 
    & \footnotesize{\emph{Assessing the Effect of Eastern Hemlock} (Tsuga canadensis) \emph{Decline from Hemlock Woolly Adelgid} (Adelges tsugae) \emph{Infestation on Ectomycorrhizal Colonization and Growth of Red Oak} (Quercus rubra) \emph{Seedlings}}\\
    \textendash ~\textsc{2020}
    & \footnotesize{\textsc{\underline{Role:} Undergraduate Research Student}} \\
    & \footnotesize{\textsc{\underline{Advisors:} Dr. Jonathan Horton, Dr. Alisa Hove}}\\
    & \footnotesize{Investigated hemlock decline from hemlock woolly adelgid by comparing ectomycorrhizal colonization rates, ectomycorrhizal taxonomic richness, and seedling growth between healthy and declining hemlock stands. Results published in \textit{American Midland Naturalist}}
    \\&\\

    \textsc{2019} 
    & \footnotesize{\emph{A CEREUS (Consortium Exchanging Research Experiences for Undergraduate Students) Approach to Assessing Ecological Responses in the Southern Appalachians to Environmental Change}}\\
    \textendash ~\textsc{2020}
    & \footnotesize{\textsc{\underline{Role:} Research Assistant}} \\
    & \footnotesize{\textsc{\underline{Principal Investigator:} Dr. Jennifer Rhode Ward}} \\
    & \footnotesize{Tracked the effects of climate variability on Southern Appalachian perennial plant phenology and monitored regrowth of plant communities following invasive species removal
}
    \\&\\

    \textsc{2018} 
    & \footnotesize{\emph{Factors Contributing to Reproductive Success and Failure in Virginia Spiraea} (Spiraea virginiana)\emph{, a Federally Threatened Shrub}}\\
    \textendash ~\textsc{2020}
    & \footnotesize{\textsc{\underline{Role:} Undergraduate Research Student}} \\
    & \footnotesize{\textsc{\underline{Advisor:} Dr. Jennifer Rhode Ward}}\\
    & \footnotesize{Examined Virginia meadowsweet (\textit{Spiraea virginiana}) reproduction, breeding, and physiological responses to light availability to assess barriers to species survival. Facilitated native plant education and contemplative practice workshops for pre-K\textendash  middle school students. Reintroduced propagule of only known living member of an extirpated \textit{S. virginiana} population, which had been maintained in the Living Collections of the Arnold Arboretum, to the Asheville Botanical Garden in the same county where the specimen was collected.
    by comparing reproductive effort and output among populations, assessed seed production by manipulating breeding systems, identified insect visitors to ascertain potential pollinators, used physiological measurements in experimental gardens to determine whether the species is shade-tolerant, and tested the potential for hybridization with Japanese spiraea (\textit{Spiraea japonica}), an invasive congener.
    
}
\end{tabular}


%See full list of projects on \website{example.com/projects}
%\vspace{1em}

%\begin{minipage}[t]{0.505\textwidth}

%\outerlist{

%\entrybig[\textbullet]
%{Project Title (Technology Used, 2019)}{}
%{Short explanation of the project}{}

%}

%\end{minipage}
%\begin{minipage}[t]{0.48\textwidth}

%\outerlist{

%\entrybig[\textbullet]
%{Project Title (Technology Used, 2019)}{}
%{Short explanation of the project}{}

%}

%\end{minipage}